\begin{table}[ht!]
\centering

% Caption 与表格顶线之间的距离
\setlength{\belowcaptionskip}{6pt}

\caption{
\textbf{Hierarchical skills extracted from on-policy trajectories.}
For each dataset, we show one successful and one failed trajectory.
Episode-level skills summarize reusable global behavior, while critical-step
skills target sparse decision points. Step indices are 0-based analyzer keys.
}
\label{tab:hierarchical-skills-success-failure}

% 仅缩小表格内容，不缩小 caption
\scriptsize
\setlength{\tabcolsep}{2.4pt}
\renewcommand{\arraystretch}{1.10}

\begin{tabularx}{\linewidth}{
  >{\raggedright\arraybackslash}p{0.10\linewidth}
  >{\raggedright\arraybackslash}p{0.07\linewidth}
  >{\raggedright\arraybackslash}p{0.18\linewidth}
  >{\raggedright\arraybackslash}p{0.28\linewidth}
  >{\raggedright\arraybackslash}X
}
\toprule
\textbf{Dataset}
&
\textbf{Outcome}
&
\textbf{Task}
&
\textbf{Episode-level skill}
&
\textbf{Critical step skills}
\\
\midrule

\textsc{ALFWorld}
&
Success
&
clean some kettle and put it in cabinet.
&
Workflow: first locate and take the target object, then move to the cleaning
station (sinkbasin) to clean it, then go to a suitable storage location
(cabinet), open it if closed, and finally place the object inside.
&
\textbf{t=0} Go directly to the countertop or likely surface where the kettle
could be. \newline
\textbf{t=2} After acquiring the kettle, immediately go to the sinkbasin to
clean it. \newline
\textbf{t=4} After cleaning, go to a cabinet (cabinet 1) rather than pausing.
\newline
\textbf{t=6} Open the closed cabinet if needed before placing the object
inside.
\\

\addlinespace[0.30em]

\textsc{ALFWorld}
&
Failure
&
put a clean soapbar in cart.
&
Avoid placing a soapbar in the cart without first confirming it is clean.
The core mistake is ignoring the cleanliness requirement; the warning sign is
repeatedly moving the soapbar without checking or cleaning it.
&
\textbf{t=1} Take and examine the soapbar to determine if it needs cleaning.
\newline
\textbf{t=2} If the soapbar is dirty, clean it using a sink or appropriate
tool before moving to the cart.
\\

\midrule

\textsc{WebShop}
&
Success
&
Find me makeup remover for sensitive skin, nail polish with style: lagom
5 layer cotton pad, and price lower than 40.00 dollars.
&
Search broadly for the specific product name and key constraints, then click
the first matching product result to view details, verify the attributes and
price, and click 'Buy Now' to finalize.
&
\textbf{t=1} First, click the most relevant product result
(the LAGOM cotton pad) to view its detailed page.
\\

\addlinespace[0.30em]

\textsc{WebShop}
&
Failure
&
Find coffee tables with steel frame, storage space, brown, size with shelf,
and price below \$110.
&
Core mistake: Buying a product without confirming it has a steel frame.
Warning signs: product title lacks mention of 'steel frame'; search results
include many unrelated items; product page shows color and size filters but
not frame material. Avoid relying on partial matches; always verify all
specific attributes, especially material, before finalizing purchase.
&
\textbf{t=1} Before clicking on a product, examine its title and check if it
explicitly mentions steel frame or other required attributes. \newline
\textbf{t=2} On the product page, click on 'Description' or 'Features' to
verify the steel frame and shelf size before clicking 'Buy Now'.
\\

\midrule

\textsc{Search}
&
Success
&
Who illustrated Hunter S. Thompson's novel Fear and Loathing in Las Vegas?
&
Workflow: First, query using core entities (author/title) to gather context;
if initial search lacks direct answer, reformulate query specifically targeting
the required attribute (illustrator) and use the new results to extract the
answer.
&
\textbf{t=1} If the initial search results do not directly answer the question,
reformulate the search query to specifically target the missing attribute
(here, 'illustrated by').
\\

\addlinespace[0.30em]

\textsc{Search}
&
Failure
&
What is the full founding date of GroenLinks, the party led by Jesse Feras
Klaver?
&
Avoid ignoring crucial temporal precision; when the task demands a specific
date, search or extract the full date, not just the year, even if the year is
initially prominent. Warning sign: Answering with only a year when documents
contain more precise information.
&
\textbf{t=2} Extract the full founding date from documents about GroenLinks,
not just the year.
\\

\bottomrule
\end{tabularx}
\end{table}